\documentclass[11pt]{article}
\usepackage[margin=0.9in]{geometry}
\usepackage{booktabs}
\usepackage{amsmath}
\usepackage[hidelinks]{hyperref}
\title{Chunk-Length Switching Outside RoboCasa:\\LIBERO and MimicGen Results}
\author{Long-horizon adaptive chunking project}
\date{July 29, 2026}
\begin{document}
\maketitle

\section*{Setup}
Same recipe as the RoboCasa study, on diffusion policies. We label
demonstration states with the open-loop stability exponent $\lambda$
(perturb the first action, replay the chunk open loop, fit the log
divergence slope), train small predictor heads on the labels (vision,
proprioception, and V-JEPA feature variants), and evaluate fixed chunk
lengths $k \in \{16, 8, 4, 1\}$ against switching $(k_{\text{stable}},
k_{\text{unstable}})$ with $(16,4)$ and $(16,1)$, plus $(8,\cdot)$
completions on the tasks where fixed 8 is competitive. Heads are queried at
every replan. 50 episodes per cell.

\section*{LIBERO}
\begin{table}[h]
\centering
\small
\setlength{\tabcolsep}{3.5pt}
\begin{tabular}{lccccccccccccc}
\toprule
& \multicolumn{4}{c}{Fixed $k$} & \multicolumn{3}{c}{vjepa} & \multicolumn{3}{c}{vision} & \multicolumn{3}{c}{proprio} \\
\cmidrule(lr){2-5}\cmidrule(lr){6-8}\cmidrule(lr){9-11}\cmidrule(lr){12-14}
Task & 16 & 8 & 4 & 1 & (16,4) & (16,1) & (8,1) & (16,4) & (16,1) & (8,1) & (16,4) & (16,1) & (8,1) \\
\midrule
Kitchen 3 moka pot   & \textbf{.72} & .44 & .22 & .08 & .36 & .40 & .20 & .56 & .42 & .24 & .30 & .10 & .10 \\
Kitchen 4 bowl in drawer & .76 & \textbf{.84} & .76 & .46 & .54 & .62 & .82 & .64 & .58 & .72 & .58 & .38 & .58 \\
Kitchen 6 mug in microwave & .00 & .00 & .00 & .00 & .00 & .00 & -- & .00 & .00 & -- & .00 & -- & -- \\
Living Room 2 soup and sauce & .00 & .00 & .00 & .00 & .00 & .00 & -- & .00 & .00 & -- & .00 & .00 & -- \\
Living Room 2 cheese and butter & .00 & .00 & .00 & .00 & .00 & .00 & -- & .00 & .00 & -- & .00 & .00 & -- \\
\bottomrule
\end{tabular}
\caption{LIBERO success rates. Bold is the best cell per task. Kitchen 6 is
zero at every completed cell (its proprio $(16,1)$ cell is still running);
dashes were not run.}
\end{table}

\section*{MimicGen}
\begin{table}[h]
\centering
\small
\setlength{\tabcolsep}{4pt}
\begin{tabular}{lcccccccccccc}
\toprule
& \multicolumn{4}{c}{Fixed $k$} & \multicolumn{4}{c}{full} & \multicolumn{2}{c}{vision} & \multicolumn{2}{c}{proprio} \\
\cmidrule(lr){2-5}\cmidrule(lr){6-9}\cmidrule(lr){10-11}\cmidrule(lr){12-13}
Task & 16 & 8 & 4 & 1 & (16,4) & (16,1) & (8,4) & (8,1) & (16,4) & (16,1) & (16,4) & (16,1) \\
\midrule
Nut assembly & .34 & .42 & .42 & \textbf{.52} & .40 & .32 & .34 & .42 & .34 & .46 & .44 & .38 \\
Coffee preparation & .18 & \textbf{.24} & .10 & .04 & .16 & .10 & .24 & .06 & .12 & .12 & .02 & .10 \\
\bottomrule
\end{tabular}
\caption{MimicGen success rates with in-domain heads. Bold is the best cell
per task.}
\end{table}

\section*{Reading}
The fixed-$k$ profiles match the label story. Kitchen 3 collapses from .72
to .08 as the chunk shrinks, the signature of a task where frequent
replanning itself destabilizes the policy. Nut assembly is the opposite:
mostly unstable, so the shortest chunk is best at .52. Coffee preparation
needs a middle chunk (.24 at $k{=}8$) and dies at $k{=}1$.

Learned switching never beats the best fixed chunk on these platforms,
though from the right base it can now match it. The $(8,\cdot)$
completions sharpen the picture: Kitchen 4's vjepa $(8,1)$ reaches .82
against the best fixed .84, and coffee's full-head $(8,4)$ reaches .24,
exactly its best fixed chunk. Nut assembly's $(8,1)$ lands on fixed 8 at
.42, still below fixed 1 at .52. The closed-loop unstable tasks pay for
every false positive: all three Kitchen 3 $(8,1)$ cells (.10 to .24) fall
far below fixed 16 at .72, and coffee $(8,1)$ collapses to .06. The
pattern is the same one we measured directly on RoboCasa: the heads rank
states well offline (AUROC .78 to .89) but the stability signal shifts
between demonstration states and rollout states, so the threshold fires in
the wrong places. On a task like Kitchen 3 the cost is largest, because
every false positive buys a step of the closed-loop instability that made
short chunks fail there in the first place; its proprio $(16,1)$ cell lands
at .10, essentially the fixed $k{=}1$ result.

In the two-axis taxonomy: Kitchen 3 and coffee sit in the closed-loop
unstable row where false positives are expensive and long or middle chunks
are mandatory. Nut assembly is mostly open-loop unstable, and the right
answer is a short fixed chunk rather than switching. Kitchen 6 and both
Living Room 2 tasks are at zero everywhere, unstable on both axes, and no
chunk length can help. These results complete the negative half of the
paper's claim: the taxonomy predicts not only where switching wins on
RoboCasa but also where it must lose elsewhere.
\end{document}
